\section{Empirical Evaluation}
\label{sec:eval}

We evaluate \rnwise{} against five baselines and across four domains, addressing
four research questions:
\begin{itemize}
\item[\textbf{RQ1}] \emph{Coverage \& fault detection.} Does \rnwise{} achieve
higher output coverage and fault detection than established baselines?
\item[\textbf{RQ2}] \emph{Generality.} Does the same algorithm transfer across
tabular ML, vision, NLP, and quantum systems?
\item[\textbf{RQ3}] \emph{Strength.} Does the method scale to higher interaction
strengths $s$, and does the $M\ge v^{s}$ bound hold?
\item[\textbf{RQ4}] \emph{Cost \& sensitivity.} What is the computational cost,
and how sensitive is it to design choices?
\end{itemize}

\subsection{Experimental setup}
The primary study uses the UCI Adult income dataset~\cite{dua2019uci} with an
XGBoost classifier~\cite{chen2016xgboost} ($\approx 87\%$ test accuracy). We
abstract its behaviour into a $q=9$, all-ternary output space (pairwise universe
$\binom{9}{2}\cdot 3^{2}=324$) spanning the decision, confidence, decision margin,
sex-flip fairness consistency, age band, employment group, education- and
hours-response monotonicity, and capital-gain regime. Eight behavioural faults are
seeded as rare-but-reachable pairwise output signatures (e.g.\ fairness
inconsistency localised to young applicants; a non-monotone education response
under a positive decision).

Baselines are scored on the identical feasible output universe: classical
input-side combinatorial testing (Input CT), random testing at matched budget,
property-based testing, metamorphic testing, and DeepCT-style
activation-combination sampling. All experiments take an explicit seed;
significance uses 30 fixed seeds. Reported values are means with $95\%$ percentile
bootstrap confidence intervals; pairwise comparisons use the one-sided
Mann--Whitney $U$ test with Cliff's $\delta$ effect size.

\subsection{RQ1: Coverage and fault detection}
Table~\ref{tab:main} reports the 30-run comparison. \rnwise{} attains
\emph{perfect} pairwise output coverage on every run ($\OCov_2 = 100\%$,
CI $[100,100]$) with $\approx 20$ prioritised tests, and $98.3\%$ fault detection.
The strongest baseline, DeepCT, reaches $91.2\%$ coverage and $99.6\%$ fault
detection but requires a large random activation pool to do so; the four classical
baselines remain between $52\%$ and $59\%$ coverage and below $35\%$ fault
detection.

\begin{table}[t]
\centering
\caption{Main comparison on UCI Adult + XGBoost (pairwise, 30 runs).
Means with $95\%$ bootstrap CIs.}
\label{tab:main}
\begin{tabular}{lrrr}
\toprule
Method & Tests & OCov$_s$ (95\% CI) & FDR (95\% CI) \\
\midrule
\textbf{Reverse N-Wise} & 20 & 100.0\% [100.0, 100.0] & 98.3\% [96.7, 99.6] \\
Input CT & 13 & 53.9\% [51.5, 56.4] & 35.0\% [29.6, 40.4] \\
Random & 20 & 52.5\% [50.4, 54.9] & 31.2\% [24.6, 37.9] \\
Property-Based & 20 & 58.6\% [56.6, 60.7] & 34.2\% [28.3, 40.0] \\
Metamorphic & 19 & 53.4\% [49.8, 57.3] & 28.7\% [22.9, 34.6] \\
DeepCT & 17 & 91.2\% [91.1, 91.3] & 99.6\% [98.8, 100.0] \\
\bottomrule
\end{tabular}
\end{table}

Table~\ref{tab:sig} gives the pairwise significance tests. \rnwise{} dominates
\emph{every} baseline on output coverage---including DeepCT---at
$p\approx 6\times10^{-13}$ with Cliff's $\delta=1.0$ (perfect, large effect):
on all 30 seeds it strictly out-covers each competitor. On fault detection it
dominates the four classical baselines ($p<10^{-12}$) and is statistically
indistinguishable from DeepCT ($p=0.92$, $\delta=-0.10$, negligible). We report
this tie transparently: both methods essentially detect all reachable seeded
faults; the advantage of \rnwise{} is that it reaches this by \emph{direct
inversion} rather than by sampling a large random pool (see RQ4).

\begin{table}[t]
\centering
\caption{Pairwise Mann--Whitney $U$ (\rnwise{} $>$ baseline) with Cliff's
$\delta$. $^{*}$ denotes $p<0.05$.}
\label{tab:sig}
\begin{tabular}{llrrl}
\toprule
Metric & vs.\ Baseline & $p$-value & $\delta$ & Effect \\
\midrule
$\mathrm{OCov}_s$ & Input CT & 6.0e-13$^{*}$ & 1.00 & large \\
$\mathrm{OCov}_s$ & Random & 6.1e-13$^{*}$ & 1.00 & large \\
$\mathrm{OCov}_s$ & Property-Based & 6.1e-13$^{*}$ & 1.00 & large \\
$\mathrm{OCov}_s$ & Metamorphic & 6.1e-13$^{*}$ & 1.00 & large \\
$\mathrm{OCov}_s$ & DeepCT & 5.8e-13$^{*}$ & 1.00 & large \\
$\eta_s$ & Input CT & 6.2e-08$^{*}$ & 0.80 & large \\
$\eta_s$ & Random & 1.5e-11$^{*}$ & 1.00 & large \\
$\eta_s$ & Property-Based & 1.5e-11$^{*}$ & 1.00 & large \\
$\eta_s$ & Metamorphic & 1.5e-11$^{*}$ & 1.00 & large \\
$\eta_s$ & DeepCT & 1.000 & -0.62 & large \\
FDR & Input CT & 1.6e-12$^{*}$ & 1.00 & large \\
FDR & Random & 1.8e-12$^{*}$ & 1.00 & large \\
FDR & Property-Based & 1.6e-12$^{*}$ & 1.00 & large \\
FDR & Metamorphic & 1.6e-12$^{*}$ & 1.00 & large \\
FDR & DeepCT & 0.920 & -0.10 & negligible \\
\bottomrule
\end{tabular}
\end{table}

\subsection{RQ2: Multi-domain generality}
Table~\ref{tab:multidomain} applies the \emph{same} four-stage algorithm---changing
only the SUT, the output partition, the search space, and the inverse-map
optimiser---to a noisy quantum circuit (variational/VQE inverse map), an MLP image
classifier on digits (over a PCA latent space), and a text classifier on
newsgroups (over a TF--IDF$+$SVD embedding space). In every domain \rnwise{}
reaches $100\%$ feasible output coverage and detects all five domain-specific
seeded faults, whereas matched-budget random testing covers only $62\%$--$74\%$
and detects $0$--$2$ of $5$. Together with the tabular study, this exercises the
method on four structurally different system types with no change to the core
algorithm.

\begin{table}[t]
\centering
\caption{Multi-domain generality (per domain: \rnwise{} vs.\ matched-budget
random).}
\label{tab:multidomain}
\begin{tabular}{llrrrr}
\toprule
Domain & Opt. & Method & Tests & OCov$_s$ & FDR \\
\midrule
Quantum & VQE & \textbf{Reverse N-Wise} & 13 & 100.0\% & 5/5 \\
 &  & Random & 13 & 73.0\% & 2/5 \\
\midrule
Vision & Jaya & \textbf{Reverse N-Wise} & 15 & 100.0\% & 5/5 \\
 &  & Random & 15 & 61.6\% & 0/5 \\
\midrule
NLP & Jaya & \textbf{Reverse N-Wise} & 17 & 100.0\% & 5/5 \\
 &  & Random & 17 & 74.1\% & 1/5 \\
\bottomrule
\end{tabular}
\end{table}

\subsection{RQ3: Interaction strength}
Table~\ref{tab:strength} varies the coverage strength $s\in\{2,3,4\}$ on the
tabular SUT. \rnwise{} maintains $100\%$ feasible coverage at every strength while
the prioritised suite grows far more slowly than the output universe
($324\to2{,}268\to10{,}206$ tuples but only $19\to58\to126$ tests). The
Corollary~\ref{cor:lb} lower bound $M\ge v^{s}$ ($9,27,81$) holds throughout,
empirically corroborating the theory.

\begin{table}[t]
\centering
\caption{Strength scaling on UCI Adult + XGBoost. The $M\ge v^{s}$ bound
(Corollary~\ref{cor:lb}) holds at every strength.}
\label{tab:strength}
\begin{tabular}{rrrrrrl}
\toprule
$s$ & Universe & Feasible & Tests & OCov$_s$ & $v^s$ & Bound \\
\midrule
2 & 324 & 257 & 19 & 100.0\% & 9 & $\checkmark$ \\
3 & 2268 & 1343 & 58 & 100.0\% & 27 & $\checkmark$ \\
4 & 10206 & 3931 & 126 & 100.0\% & 81 & $\checkmark$ \\
\bottomrule
\end{tabular}
\end{table}

\subsection{RQ4: Cost and sensitivity}
Table~\ref{tab:cost} profiles the pipeline on commodity hardware. The end-to-end
run completes in under $10$ seconds using $\approx 68$ MB, and issues about
$8{,}200$ black-box SUT evaluations---roughly $8\%$ of a modest exhaustive-input
sweep. The feasibility scan dominates the wall-clock time; warm-started inverse
mapping is inexpensive.

\begin{table}[t]
\centering
\caption{Cost profile of the \rnwise{} pipeline (UCI Adult + XGBoost).}
\label{tab:cost}
\begin{tabular}{lr}
\toprule
Quantity & Value \\
\midrule
Stage 1 (feasibility) & 8.74 s \\
Stage 2 (covering array) & 0.00 s \\
Stage 3 (inverse mapping) & 0.72 s \\
Stage 4 (prioritise) & 0.23 s \\
Pipeline total & 9.68 s \\
Peak memory & 68.5 MB \\
SUT evaluations & 8{,}217 \\
Fraction of exhaustive sweep & 8.22\% \\
\bottomrule
\end{tabular}
\end{table}

Ablations (omitted here for space, included in the artifact) confirm the design
choices: (i) among inverse-map optimisers, Jaya and Whale are the most efficient
(comparable realisation at $\approx 7\times$ fewer evaluations than Firefly),
justifying Jaya as the default; (ii) coverage remains $100\%$ as output
granularity varies from $2$ to $4$ bins per continuous factor; and (iii) a mild
input regulariser ($\lambda>0$) slightly improves realisation by keeping
synthesised inputs near the data manifold.
